% Esquisse d'un programme (cote 119, copy 1): archive cote 119, site page 15 (PDF page 16), typescript page 13.
% typescript. Draft compared: EsquisseFr.pdf (text layer):490-541.
% AI-transcribed by Claude Opus 5.5 (claude-opus-5-5, adaptive thinking, effort high), Message Batches API, batch esq1 (msgbatch_01LNsogFyF3y6KSp9UbtJtBk), 2026-11-12; not yet checked by eye.
\scanpage{119/15}
\mspage{13}

nant nécessairement naissance à des courbes de genre zéro (qui ne
fournissent donc pas des ``modules''), on peut dire que la courbe
en question est ``épinglée'' dès qu'on fixe trois de ses points, par
exemple trois sommets de la carte, ou plus généralement trois cen-%
tres de facettes (sommets, arêtes ou faces) - dès lors l'applica-%
tion structurale $f : X \to \hand{$\mathbb{P}$}^{1}_{\hand{$\mathbb{C}$}}$ peut s'interpréter comme une fonction
rationnelle
\[ f(z) = P(z)/Q(z) \hand{$\in$} \hand{$\mathbb{C}$}(z) \]
bien déterminée, quotient de deux polynômes bien déterminés pre-%
miers entre eux avec $Q$ unitaire, satisfaisant à des conditions
algébriques qui traduisent notamment le fait que $f$ soit non
ramifié en dehors des valeurs $0$, $1$, $\infty$, et qui impliquent que les
coefficients de ces polynômes sont des \uline{nombres algébriques} ; donc
leurs zéros sont des nombres algébriques, qui représentent res-%
pectivement les sommets et les centres des faces de la carte en-%
visagée.

Revenant au cas général, les cartes finies s'interpré-%
tant comme des revêtements sur $\bar{\hand{$\mathbb{Q}$}}$ d'une courbe algébrique définie
sur le corps premier \hand{$\mathbb{Q}$} lui-même, il en résulte que le groupe de
Galois \hand{$\mathbb{I}\Gamma$} de $\bar{\hand{$\mathbb{Q}$}}$ sur \hand{$\mathbb{Q}$} opère sur la catégorie de ces cartes de
façon naturelle. Par exemple, l'opération d'un automorphisme
$\gamma \hand{$\in$} \hand{$\mathbb{I}\Gamma$}$ sur une carte sphérique donnée par la fonction rationnelle
ci-dessus, est obtenue en appliquant $\gamma$ aux coefficients des po-%
lynômes $P$, $Q$. Voici donc ce mystérieux groupe \hand{$\mathbb{I}\Gamma$} intervenir comme
agent transformateur sur des formes topologico-combinatoires de la
nature la plus élémentaire qui soit, amenant à se poser des ques-%
tions comme : telles cartes orientées données sont-elles ``conjuguées'',
ou : quels exactement sont les conjuguées de telle carte orientée
donnée ? (Il y en a, visiblement, un nombre fini seulement.)

J'ai traité quelques cas concrets (pour des revêtements
de bas degrés) par des expédients divers, J. Malgoire en a traité
quelques autres - je doute qu'il y ait une méthode uniforme permet-%
tant d'y répondre à coups d'ordinateurs. Ma réflexion très vite

\margin{15}
